\chapter{Glossary}
\label{app:glossary}

\begin{description}[leftmargin=3.2cm,style=nextline]
\item[AGC] Automatic gain control; a receiver feedback loop that
reduces gain for strong signals and restores it for weak ones.
\item[AGN, BTU, FB, \dots] see CW operating abbreviations,
Chapter~\ref{ch:phonetics-codes}.
\item[AM] Amplitude modulation.
\item[APRS] Automatic Packet Reporting System; packet-radio-based
position and status broadcasting network.
\item[ARQ] Automatic repeat request; retransmission triggered by a
failed error-detection check.
\item[Balun] Balanced-to-unbalanced transformer.
\item[Bandwidth] The range of frequencies a signal or filter occupies.
\item[BJT] Bipolar junction transistor.
\item[CEPT] European Conference of Postal and Telecommunications
Administrations; facilitates reciprocal amateur licensing in Europe.
\item[Critical frequency] Highest frequency reflected by the ionosphere
for a vertically transmitted signal.
\item[CRC] Cyclic redundancy check; an error-detection scheme.
\item[CTCSS] Continuous tone-coded squelch system; sub-audible repeater
access tone.
\item[CW] Continuous wave; on-off keyed Morse code telegraphy.
\item[dB, dBm, dBi, dBd] Decibel and its absolute-reference variants;
see Chapter~\ref{ch:decibels}.
\item[Deviation] Peak frequency shift in FM, $\Delta f$.
\item[Dipole] A two-element antenna fed at its center.
\item[Directivity] Concentration of an antenna's radiation in
particular directions.
\item[Duplex] Simultaneous two-way communication (full duplex) or
alternating two-way communication (half-duplex).
\item[EMF] Electromotive force; the source-side driving potential of a
circuit.
\item[FEC] Forward error correction.
\item[FM] Frequency modulation.
\item[FSK] Frequency-shift keying.
\item[Greyline] The twilight band circling Earth between day and night,
favorable for certain HF propagation.
\item[IARU] International Amateur Radio Union.
\item[Impedance] The total (resistive plus reactive) opposition to AC
current flow, $Z = R \pm jX$.
\item[Ionosphere] The ionized upper-atmosphere region responsible for
HF skywave propagation.
\item[ITU] International Telecommunication Union.
\item[LO] Local oscillator.
\item[LUF] Lowest usable frequency for a given ionospheric path.
\item[MOSFET] Metal-oxide-semiconductor field-effect transistor.
\item[MUF] Maximum usable frequency for a given ionospheric path.
\item[Op-amp] Operational amplifier.
\item[PM] Phase modulation.
\item[PSK] Phase-shift keying.
\item[Q (quality factor)] Ratio of energy stored to energy dissipated
per cycle in a resonant circuit; sets resonance sharpness.
\item[QAM] Quadrature amplitude modulation.
\item[Q-code] Standardized three-letter operating shorthand beginning
with Q (QTH, QSL, QRM, \dots); see Chapter~\ref{ch:phonetics-codes}.
\item[Reactance] Frequency-dependent opposition to AC current from a
capacitor ($X_C$) or inductor ($X_L$), without energy dissipation.
\item[Repeater] A station that receives on one frequency and
retransmits on another to extend range.
\item[Resonance] The frequency at which inductive and capacitive
reactance cancel, $f_0 = 1/(2\pi\sqrt{LC})$.
\item[RF] Radio frequency.
\item[RMS] Root-mean-square; the AC amplitude measure equivalent to DC
heating effect.
\item[RST] Readability/Strength/Tone signal report system.
\item[RTTY] Radioteletype.
\item[Sensitivity] The weakest signal a receiver can usefully detect.
\item[Simplex] Two stations sharing one frequency, transmitting in
turn.
\item[Skin effect] The tendency of high-frequency current to crowd
toward a conductor's surface, raising effective AC resistance.
\item[Skywave] Propagation via reflection/refraction from the
ionosphere.
\item[Sporadic-E] Anomalous, patchy E-layer ionization supporting
unusual VHF/upper-HF propagation.
\item[SSB] Single-sideband modulation.
\item[SWR] Standing wave ratio.
\item[Transceiver] A combined transmitter and receiver.
\item[Yagi--Uda] A directional antenna using a driven element plus
parasitic reflector and director elements.
\end{description}
